\markboth{未来研究计划}{未来研究计划}
\chapter{未来研究计划}
\label{sec:future_intro}

\section{总结与展望}

通过工作一、工作二和工作三，本文设计并实现了一套基于\LaTeX 的上海创智学院博士开题报告撰写模板，解决了研究生在撰写开题报告过程中遇到的格式繁琐、排版困难等问题。目前，模板已经具备了基本的文档结构、样式定制和自动化编译功能，能够满足大多数用户的需求。

然而，本模板仍存在一些不足之处，需要在未来的工作中进一步完善：
\begin{enumerate}
    \item 功能扩展：目前模板主要针对开题报告，未来计划扩展支持学位论文、中期报告等更多类型的文档。
    \item 画质提升：目前的封面Logo清晰度较低，未来计划寻找更高分辨率的图像资源。
\end{enumerate}

\section{研究进度安排}

\begin{table}[H]
  \centering
  \caption{研究进度安排表}
  \begin{tabular}{|p{3cm}|p{7cm}|p{5cm}|}
  \hline
  \multicolumn{1}{|c|}{起讫时间} & \multicolumn{1}{c|}{主要研究内容} & \multicolumn{1}{c|}{预期结果、阶段性成果} \\ \hline
  2023.09 - 2024.05 & 完成工作一。 & 发表论文一。 \\[3ex] \hline
  2024.05 - 2025.05 & 完成工作二。 & 发表论文二。 \\[3ex] \hline
  2025.05 - 2026.05 & 完成工作三。 & 预期发表论文三。 \\[3ex] \hline
  2026.05 - 2027.05 & 完成工作二的扩展。 & 预期发表扩展论文一篇。 \\[3ex] \hline
  2027.05 - 2028.05 & 完成工作三的扩展。 & 预期发表扩展论文一篇。 \\[3ex] \hline
  \end{tabular}
\end{table}
